\addcontentsline{toc}{chapter}{R\'esum\'e }

Un m\'ecanisme fondamental de la cognition, n\'ecessaire \`a l'ex\'ecution de tâches de raisonnement complexes, est la capacit\'e de traiter s\'electivement les informations (attention) et de les conserver dans un \'etat accessible (m\'emoire). Nous analysons syst\'ematiquement le rôle de ces deux composantes, en commençant par l'auto-attention bas\'ee sur le modèle d'attention le plus populaire: Transformer, et en \'etendant ensuite l'architecture \`a la m\'emoire. Transformer est aujourd'hui la dernière classe d'architecture neuronale et est au coeur des demonstrations les plus fascinante du Deep Learning, il a apport\'e un changement de paradigme dans le domaine de l'intelligence artificielle. Il a remplac\'e les r\'eseaux de r\'ecurrence et de convolution par l'auto-attention comme choix architectural de facto pour la plupart des applications de l'IA.


Nous \'etudions d'abord les m\'ecanismes de calcul impliqu\'es dans un test de raisonnement visuel synth\'etique (SVRT), en analysant la capacit\'e d'une architecture de vision par ordinateur populaire (ResNet) de diff\'erentes profondeurs et entraîn\'ee sur des ensembles de donn\'ees de diff\'erentes tailles. Cela a conduit \`a une nouvelle taxonomie plus fine pour les vingt-trois tâches de SVRT, coh\'erente avec les classes de tâches de raisonnement - identiques-diff\'erentes (SD) et de relations spatiales (SR) - largement accept\'ees dans la litt\'erature. Ensuite, nous \'etudions le rôle de l'auto-attention incorpor\'ee \`a ResNet50 dans la r\'esolution du d\'efi SVRT. Inspir\'es par les deux types de systèmes d'attention visuelle, nous avons mod\'elis\'e l'auto-attention pour qu'elle soit utilis\'ee comme une attention bas\'ee sur les caract\'eristiques et sur une attention spatiale pour enrichir les cartes de caract\'eristiques d'un r\'eseau feedforward. Nous avons \'evalu\'e la capacit\'e de ces r\'eseaux d'attention \`a r\'esoudre le d\'efi SVRT et avons constat\'e que les architectures r\'esultantes \'etaient beaucoup plus efficaces pour r\'esoudre la plus difficile de ces tâches de raisonnement visuel. La nouvelle taxonomie obtenue pr\'ec\'edemment s'explique aussi partiellement par l'am\'elioration relative des deux r\'eseaux d'attention et conduit \`a des pr\'edictions testables concernant les besoins attentionnels des tâches SVRT. 

Enfin, nous d\'eveloppons une nouvelle architecture cognitive int\'egrant l'auto-attention et la m\'emoire. Nous proposons GAMR: \textbf{G}uided \textbf{A}ttention \textbf{M}odel for visual \textbf{R}easoning, motiv\'e par la th\'eorie de la vision active. Le GAMR a des m\'ecanismes de fonctionnement similaires \`a ceux du cerveau qui r\'esout des tâches complexes de raisonnement visuel par des s\'equences de changements d'attention pour s\'electionner et acheminer en m\'emoire les informations visuelles pertinentes pour la tâche. Ce changement d'attention est mis en œuvre \`a l'aide d'un module d'auto-attention guid\'e par une requête g\'en\'er\'ee en interne. Nous d\'emontrons que \textit{GAMR} est efficace, robuste et compositionnel par rapport \`a l'une ou l'autre des architectures bas\'ees sur le feedforward, l'attention ou la m\'emoire. De plus, GAMR est capable de g\'en\'eraliser \`a des tâches de raisonnement complètement nouvelles. Dans l'ensemble, notre travail analyse le rôle de l'auto-attention dans l'architecture cognitive et de vision par ordinateur par leur capacit\'e \'a r\'esoudre des tâches complexes de raisonnement visuel n\'ecessitant de l'attention comme \'el\'ement cl\'e pour r\'esoudre efficacement les tâches de raisonnement.